\documentclass[a4paper,11pt]{article}
\usepackage{amsmath,amssymb,amsthm}
\usepackage[margin=2.5cm]{geometry}
\usepackage{bm}
\usepackage{algorithm}
\usepackage{algpseudocode}
\usepackage{booktabs}

\newcommand{\va}{\bm{a}}
\newcommand{\vb}{\bm{b}}
\newcommand{\vc}{\bm{c}}
\newcommand{\vct}{\tilde{\bm{c}}}
\newcommand{\ve}{\bm{e}}
\newcommand{\vu}{\bm{u}}
\newcommand{\vv}{\bm{v}}
\newcommand{\Xb}[2]{X_{#1}^{#2}}
\newcommand{\Xn}[2]{\bar{X}_{#1}^{#2}}
\newcommand{\Gb}[1]{G_{#1}}
\newcommand{\frob}[1]{\lVert #1 \rVert_{\mathrm{F}}}
\newcommand{\tr}{\mathrm{tr}}
\newcommand{\Nk}{N_k}
\newcommand{\Nt}{N_\tau}

\theoremstyle{plain}
\newtheorem{theorem}{Theorem}

\title{Canonical Correlation Analysis for Joint Eigenfunctions\\of Cosmological Perturbation Bases}
\author{Will Handley \and Wei Ning Deng}
\date{\today}

\begin{document}
\maketitle

\begin{abstract}
We formulate the problem of finding joint eigenfunctions of two cosmological perturbation bases as canonical correlation analysis (CCA) in $k$-space.
This provides a principled framework for the stacked-SVD approach, makes the choice of weighting across perturbation types explicit, and clarifies the conditions under which the results are independent of this choice.
We present the complete pipeline from raw Boltzmann code output to a sparse set of joint eigenfunctions.
\end{abstract}

%----------------------------------------------------------------------
\section{Problem Statement}
\label{sec:problem}
%----------------------------------------------------------------------

A Boltzmann code produces solutions for cosmological perturbation variables $\delta_r, \delta_m, v_r, v_m$ (and potentially others) as functions of conformal time $\tau$ and wavenumber $k$.
For a given set of $\Nk$ wavenumbers, we have two bases of solutions:
\begin{itemize}
    \item \textbf{Basis 1} (e.g.\ integer-$k$): matrices $\Xb{1}{p}$ of shape $(\Nt, \Nk)$ for each perturbation type $p$,
    \item \textbf{Basis 2} (e.g.\ allowed-$k$): matrices $\Xb{2}{p}$ of the same shape,
\end{itemize}
where rows index conformal time and columns index $k$-modes.

We seek coefficient vectors such that the time evolution produced by basis~1 matches that of basis~2, for all perturbation types simultaneously.
The final result should be a set of \emph{sparse} coefficient vectors, each concentrated on one or a few $k$-modes, representing the joint eigenfunctions of the two bases.

%----------------------------------------------------------------------
\section{The Stacked-SVD Approach}
\label{sec:stacking}
%----------------------------------------------------------------------

The current method normalises each perturbation type per basis by its Frobenius norm $\frob{A} = \sqrt{\sum_{ij} A_{ij}^2}$, defining
\begin{equation}
    \Xn{\alpha}{p} = \frac{\Xb{\alpha}{p}}{\frob{\Xb{\alpha}{p}}},
    \label{eq:normalise}
\end{equation}
then stacks vertically:
\begin{equation}
    \bar{X}_\alpha = \begin{pmatrix} \Xn{\alpha}{\delta_r} \\ \Xn{\alpha}{\delta_m} \\ \Xn{\alpha}{v_r} \\ \Xn{\alpha}{v_m} \end{pmatrix}, \qquad \alpha \in \{1, 2\},
    \label{eq:stacking}
\end{equation}
producing matrices of shape $(4\Nt, \Nk)$, and then applies SVD to each.
While effective, this obscures the structure of the problem: the stacking is an ad hoc device to reduce a multi-field problem to a single-matrix one, and the Frobenius normalisation is a specific choice for handling the dimensional mismatch between perturbation types.

%----------------------------------------------------------------------
\section{CCA Formulation}
\label{sec:cca}
%----------------------------------------------------------------------

\subsection{Gram matrices in $k$-space}
\label{sec:gram}

For a set of positive dimensionless weights $\{w_p\}$ indexed by perturbation type, define the $k$-space Gram matrices from the normalised data~\eqref{eq:normalise}:
\begin{align}
    \Gb{1} &= \sum_p w_p \, (\Xn{1}{p})^\top \Xn{1}{p}, \label{eq:G1} \\
    \Gb{2} &= \sum_p w_p \, (\Xn{2}{p})^\top \Xn{2}{p}, \label{eq:G2} \\
    \Gb{12} &= \sum_p w_p \, (\Xn{1}{p})^\top \Xn{2}{p}, \label{eq:G12}
\end{align}
each of shape $(\Nk, \Nk)$.
Because the data is pre-normalised~\eqref{eq:normalise}, all Gram matrices are dimensionless regardless of the physical dimensions of the perturbation types.
The normalisation absorbs the per-basis, per-type scales, so that $\Gb{12}$ naturally carries the correct cross-normalisation $(\frob{\Xb{1}{p}} \frob{\Xb{2}{p}})^{-1}$ for each type.

The matrix $\Gb{\alpha}$ encodes the inner product between $k$-modes as measured across all perturbation types in basis $\alpha$, while $\Gb{12}$ encodes the cross-correlation between the two bases.

\subsection{Two-vector canonical correlation analysis}
\label{sec:cca-def}

CCA seeks two coefficient vectors $\va \in \mathbb{R}^{\Nk}$ and $\vb \in \mathbb{R}^{\Nk}$ --- one for each basis --- that maximise the correlation between the resulting time evolutions:
\begin{equation}
    \rho(\va, \vb) = \frac{\va^\top \Gb{12} \, \vb}{\sqrt{\va^\top \Gb{1} \, \va \;\cdot\; \vb^\top \Gb{2} \, \vb}}.
    \label{eq:cca}
\end{equation}
The maximisers are found by computing the SVD
\begin{equation}
    \Gb{1}^{-1/2} \, \Gb{12} \, \Gb{2}^{-1/2} = U \Sigma V^\top,
    \label{eq:svd}
\end{equation}
where $\Sigma = \mathrm{diag}(\sigma_1, \sigma_2, \ldots)$ with $\sigma_1 \geq \sigma_2 \geq \cdots \geq 0$.
The canonical correlations are $\rho_i = \sigma_i$, with corresponding coefficient vectors
\begin{align}
    \va_i &= \Gb{1}^{-1/2} \, \vu_i, \label{eq:a} \\
    \vb_i &= \Gb{2}^{-1/2} \, \vv_i, \label{eq:b}
\end{align}
where $\vu_i$ and $\vv_i$ are the $i$-th columns of $U$ and $V$ respectively.
These are already in $k$-space --- no additional back-transform is required.

Equivalently, the left canonical vectors $\va_i$ can be found from the eigenvalue problem
\begin{equation}
    \Gb{1}^{-1/2} \, \Gb{12} \, \Gb{2}^{-1} \, \Gb{12}^\top \, \Gb{1}^{-1/2} \, \vu = \lambda \, \vu,
    \label{eq:gev}
\end{equation}
with $\lambda_i = \sigma_i^2 = \rho_i^2$.
Modes with $\rho_i \approx 1$ are joint eigenfunctions; modes with $\rho_i \ll 1$ are incompatible between the two bases.

\subsection{Why whitening is necessary}
\label{sec:whitening}

A natural simplification would be to apply SVD directly to the cross-Gram matrix,
\begin{equation}
    \Gb{12} = U\Sigma V^\top,
    \label{eq:svd-G12-simple}
\end{equation}
and set $\va_i = \vu_i$, $\vb_i = \vv_i$, avoiding the need to compute $\Gb{1}^{-1/2}$
and $\Gb{2}^{-1/2}$.
This is simpler, but it solves the wrong problem.

\paragraph{Covariance versus correlation.}
The SVD of $\Gb{12}$ maximises the \emph{cross-covariance}
$\va^\top \Gb{12}\vb$ subject to Euclidean constraints $\|\va\|=\|\vb\|=1$.
CCA instead maximises the \emph{correlation}~\eqref{eq:cca}, which normalises by the
self-energy of each basis independently via $\va^\top \Gb{1}\va$ and $\vb^\top \Gb{2}\vb$.
Without this normalisation, $k$-modes that happen to carry large time-evolution amplitude
dominate the singular value spectrum even if they agree poorly between the two bases.
Furthermore, the raw singular values of $\Gb{12}$ have no natural bound and no
interpretable threshold, whereas the canonical correlations satisfy $\sigma_i\in[0,1]$
with $\sigma_i=1$ signalling an exact joint eigenfunction.

\paragraph{Frobenius normalisation is not sufficient.}
One might expect the Frobenius pre-normalisation~\eqref{eq:normalise} to already remove
this amplitude bias.
It does not, because Frobenius normalisation and $\Gb{\alpha}^{-1/2}$ whitening act at
different levels:
\begin{itemize}
    \item \textbf{Frobenius normalisation} sets $\|\Xn{\alpha}{p}\|_{\mathrm{F}}=1$ for
    each perturbation type~$p$ separately.
    It removes \emph{inter-field} amplitude differences, ensuring that
    $\delta_r,\delta_m,v_r,v_m$ contribute equally to the Gram matrices regardless of
    their physical dimensions.
    It is a property of the whole matrix, not of individual $k$-mode columns.
    \item \textbf{$\Gb{\alpha}^{-1/2}$ whitening} normalises individual \emph{$k$-mode}
    amplitudes within the combined, multi-field space.
    The diagonal entry $(\Gb{\alpha})_{ii}=\sum_p w_p\|\Xn{\alpha}{p}[\,\cdot\,,i]\|^2$
    is the total time-evolution power of mode~$i$ across all perturbation types.
    Even after Frobenius normalisation, individual modes can carry very different power,
    since $\sum_i(\Gb{\alpha})_{ii}=\sum_p w_p$ is fixed but individual entries are not.
\end{itemize}

\paragraph{Counterexample.}
Consider a single field after Frobenius normalisation, with two $k$-modes whose
column norms satisfy $\|\Xn{1}{}[\,\cdot\,,1]\|=0.01$ and
$\|\Xn{1}{}[\,\cdot\,,2]\|\approx 0.99$.
Suppose mode~1 is a \emph{perfect} joint eigenfunction ($\rho_1=1$) while mode~2
agrees only moderately ($\rho_2=0.5$).
The cross-Gram entries are then
\begin{align}
    \Gb{12}[1,1] &= 0.01^2 = 10^{-4} \quad\text{(tiny)},\\
    \Gb{12}[2,2] &\approx 0.5\times 0.99^2 \approx 0.49 \quad\text{(large)}.
\end{align}
The SVD of $\Gb{12}$ ranks mode~2 first, incorrectly.
CCA recovers the correct ordering via whitening:
\begin{equation}
    \rho_1 = \frac{\Gb{12}[1,1]}{\sqrt{\Gb{1}[1,1]\,\Gb{2}[1,1]}}
           = \frac{10^{-4}}{10^{-4}} = 1. \qquad \checkmark
\end{equation}
Frobenius normalisation set the overall matrix norm to unity, but did nothing to equalise
the column norms of the two modes; the $\Gb{\alpha}^{-1/2}$ whitening supplies the missing
per-mode normalisation.

\subsection{Relationship to stacking}
\label{sec:equivalence}

The stacked-SVD approach is equivalent to two-vector CCA on the normalised data~\eqref{eq:normalise} with unit weights $w_p = 1$.
The stacked Gram matrices are
\begin{equation}
    \bar{X}_1^\top \bar{X}_1 = \sum_p (\Xn{1}{p})^\top \Xn{1}{p} = \Gb{1}\big|_{w_p = 1},
\end{equation}
and similarly for $\Gb{2}$ and $\Gb{12}$.
The CCA formulation generalises this by allowing $w_p \neq 1$, weighting perturbation types unequally.

%----------------------------------------------------------------------
\section{Weight-Independence of Exact Matches}
\label{sec:weight-independence}
%----------------------------------------------------------------------

\begin{theorem}
\label{thm:weight-indep}
If\/ $\Xb{1}{p} \vc = \Xb{2}{p} \vc$ for all perturbation types $p$, and $\Xb{1}{p}\vc \neq 0$ for at least one $p$, then $(\va, \vb) = (\vc, \vc)$ achieves $\rho(\vc, \vc) = 1$ for any choice of positive weights $\{w_p\}$.
\end{theorem}

\begin{proof}
The hypothesis $\Xb{1}{p}\vc = \Xb{2}{p}\vc$ implies $\lVert\Xb{1}{p}\vc\rVert = \lVert\Xb{2}{p}\vc\rVert$ for each $p$, and therefore
\begin{equation}
    (\Xn{1}{p}\vc)^\top (\Xn{2}{p}\vc) = \frac{(\Xb{1}{p}\vc)^\top(\Xb{2}{p}\vc)}{\frob{\Xb{1}{p}}\frob{\Xb{2}{p}}} = \frac{\lVert\Xb{1}{p}\vc\rVert^2}{\frob{\Xb{1}{p}}\frob{\Xb{2}{p}}}.
\end{equation}
Similarly,
\begin{equation}
    \lVert\Xn{1}{p}\vc\rVert^2 = \frac{\lVert\Xb{1}{p}\vc\rVert^2}{\frob{\Xb{1}{p}}^2}, \qquad
    \lVert\Xn{2}{p}\vc\rVert^2 = \frac{\lVert\Xb{2}{p}\vc\rVert^2}{\frob{\Xb{2}{p}}^2} = \frac{\lVert\Xb{1}{p}\vc\rVert^2}{\frob{\Xb{2}{p}}^2}.
\end{equation}
Writing $s_p = \lVert\Xb{1}{p}\vc\rVert^2$, $f_p = \frob{\Xb{1}{p}}$, $g_p = \frob{\Xb{2}{p}}$:
\begin{align}
    \vc^\top \Gb{12} \, \vc &= \sum_p w_p \frac{s_p}{f_p g_p}, \\
    \vc^\top \Gb{1} \, \vc &= \sum_p w_p \frac{s_p}{f_p^2}, \\
    \vc^\top \Gb{2} \, \vc &= \sum_p w_p \frac{s_p}{g_p^2}.
\end{align}
By the Cauchy--Schwarz inequality applied to the sequences $(\sqrt{w_p s_p}/f_p)$ and $(\sqrt{w_p s_p}/g_p)$:
\begin{equation}
    \left(\sum_p w_p \frac{s_p}{f_p g_p}\right)^2 \leq \left(\sum_p w_p \frac{s_p}{f_p^2}\right)\left(\sum_p w_p \frac{s_p}{g_p^2}\right),
\end{equation}
so $\rho(\vc,\vc) \leq 1$.
Equality holds if and only if $f_p / g_p$ is the same for all $p$ with $s_p > 0$, i.e.\ $\frob{\Xb{1}{p}}/\frob{\Xb{2}{p}}$ is constant across active perturbation types.
\end{proof}

\noindent\textbf{Remark.}
In general, $\rho(\vc,\vc) = 1$ requires not only $\Xb{1}{p}\vc = \Xb{2}{p}\vc$ for all $p$, but also that the ratio $\frob{\Xb{1}{p}}/\frob{\Xb{2}{p}}$ is the same across perturbation types.
This ratio is a property of the bases, not of $\vc$.
When the two bases have similar overall amplitudes (as is typical), $f_p \approx g_p$ and $\rho(\vc,\vc) \approx 1$.

However, the two-vector CCA will still identify these modes.
If $\Xb{1}{p}\vc = \Xb{2}{p}\vc$ for all $p$, then $\Xn{1}{p}\vc$ and $\Xn{2}{p}\vc$ differ only by the ratio $g_p/f_p$, which is independent of $\vc$.
The CCA optimisation is free to choose $\va \neq \vb$ and will find $\rho = 1$ with $\va$ and $\vb$ related by a rescaling that compensates the norm ratios.
In the limit $f_p = g_p$ for all $p$, the two coefficient vectors coincide: $\va = \vb = \vc$.

%----------------------------------------------------------------------
\section{Sparse Basis via Oblique Rotation}
\label{sec:sparsity}
%----------------------------------------------------------------------

The CCA coefficient vectors $\va_i$ are not generally sparse: they are global linear combinations of $k$-modes, chosen to maximise correlation rather than sparsity.
We seek an invertible matrix $M \in \mathbb{R}^{m \times m}$ such that the rotated vectors
\begin{equation}
    \vct_j = \sum_{i=1}^{m} M_{ji} \, \va_i, \qquad \text{i.e.} \quad \tilde{A} = A M^\top, \quad A = [\va_1 \mid \cdots \mid \va_m] \in \mathbb{R}^{\Nk \times m},
    \label{eq:rotation}
\end{equation}
are sparse (concentrated on one or a few $k$-modes).
Orthogonality of the $\vct_j$ is not required --- only sparsity.

\subsection{Why Varimax fails}
\label{sec:varimax}

Varimax rotation requires orthonormal input.
The CCA coefficient vectors $\va_i$ are not orthonormal in $k$-space in general.
Applying Varimax directly would produce incorrect results.

\subsection{Oblique rotation (Promax)}
\label{sec:promax}

Promax generalises Varimax by relaxing the orthogonality constraint.
It finds $M$ such that the columns of $\tilde{A} = A M^\top$ have ``simple structure'' (maximally sparse), with $M$ invertible but not necessarily orthogonal.
The additional degrees of freedom allow sparser solutions than any orthogonal rotation.

\subsection{Independent component analysis (FastICA)}
\label{sec:ica}

An alternative approach models the problem as blind source separation: the sparse $\vct_j$ are ``independent sources'' mixed by a matrix to produce the observed $\va_i$.
FastICA recovers the sources by maximising non-Gaussianity (kurtosis), which for peaked distributions is equivalent to maximising sparsity.

\subsection{Comparison of methods}
\label{sec:comparison}

\begin{center}
\begin{tabular}{llll}
\toprule
Method & Orthogonality & Sparsity objective & Notes \\
\midrule
Varimax & Required & Maximise kurtosis of columns & Not applicable \\
\textbf{Promax} & \textbf{Not required} & \textbf{Simple structure / kurtosis} & \textbf{Recommended} \\
FastICA & Not required & Maximise non-Gaussianity & Excellent alternative \\
$\ell_1$ + log-det & Not required & Direct $\ell_1$ minimisation & Most rigorous \\
\bottomrule
\end{tabular}
\end{center}

%----------------------------------------------------------------------
\section{Complete Pipeline}
\label{sec:pipeline}
%----------------------------------------------------------------------

\begin{algorithm}[H]
\caption{CCA-based sparse joint eigenfunction method}
\label{alg:pipeline}
\begin{algorithmic}[1]
\Statex \textbf{Input:} $\Xb{1}{p}, \Xb{2}{p} \in \mathbb{R}^{\Nt \times \Nk}$ for each perturbation type $p$; weights $\{w_p\}$; correlation threshold $\rho_{\min}$
\Statex

\State \textbf{Normalise} \Comment{\S\ref{sec:gram}}
\State \quad $\Xn{\alpha}{p} \gets \Xb{\alpha}{p} / \frob{\Xb{\alpha}{p}}$ for all $\alpha, p$
\Statex

\State \textbf{Compute Gram matrices} \Comment{\S\ref{sec:gram}}
\State \quad $\Gb{1} \gets \sum_p w_p \, (\Xn{1}{p})^\top \Xn{1}{p}$ \Comment{$(\Nk, \Nk)$}
\State \quad $\Gb{2} \gets \sum_p w_p \, (\Xn{2}{p})^\top \Xn{2}{p}$
\State \quad $\Gb{12} \gets \sum_p w_p \, (\Xn{1}{p})^\top \Xn{2}{p}$
\Statex

\State \textbf{Solve CCA} \Comment{\S\ref{sec:cca-def}}
\State \quad $U \Sigma V^\top \gets \mathrm{SVD}(\Gb{1}^{-1/2} \, \Gb{12} \, \Gb{2}^{-1/2})$
\State \quad Keep $m$ components with $\sigma_i > \rho_{\min}$
\State \quad $\va_i \gets \Gb{1}^{-1/2} \vu_i$ for $i = 1, \ldots, m$ \Comment{basis-1 $k$-space coefficients}
\State \quad $\vb_i \gets \Gb{2}^{-1/2} \vv_i$ for $i = 1, \ldots, m$ \Comment{basis-2 $k$-space coefficients}
\Statex

\State \textbf{Oblique rotation for sparsity} \Comment{\S\ref{sec:sparsity}}
\State \quad $A \gets [\va_1 \mid \cdots \mid \va_m]$ \Comment{$(\Nk, m)$}
\State \quad $\tilde{A} \gets \mathrm{Promax}(A)$ \Comment{or FastICA}
\Statex

\Statex \textbf{Output:} $\tilde{A} \in \mathbb{R}^{\Nk \times m}$
\Statex \quad Each column: sparse coefficient vector for a valid joint eigenfunction
\Statex \quad For modes with $\rho \approx 1$: $\va_i \approx \vb_i$ (same coefficients in both bases)
\end{algorithmic}
\end{algorithm}

%----------------------------------------------------------------------
\section{Discussion}
\label{sec:discussion}
%----------------------------------------------------------------------

\begin{enumerate}
    \item \textbf{Principled framework.} CCA is a well-studied statistical method with known optimality properties. The stacked-SVD approach is recovered as the special case $w_p = 1$ on the normalised data $\Xn{\alpha}{p}$.

    \item \textbf{Explicit weighting.} The dependence on perturbation-type weights $\{w_p\}$ is made transparent. Pre-normalisation~\eqref{eq:normalise} renders the data dimensionless, so the weights $w_p$ are dimensionless and encode only relative importance.

    \item \textbf{Works in $k$-space.} The Gram matrices are $(\Nk, \Nk)$, compared to $(4\Nt, \Nk)$ for the stacked matrices. For $\Nt \gg \Nk$ this is more efficient.

    \item \textbf{No back-transform.} The CCA coefficient vectors~\eqref{eq:a}--\eqref{eq:b} are already in $k$-space. The stacked-SVD approach requires an explicit back-transform through the right singular vectors.

    \item \textbf{Two coefficient vectors.} CCA naturally produces separate coefficient vectors $\va_i$ and $\vb_i$ for each basis. For well-matched modes ($\rho_i \approx 1$), these nearly coincide. The stacked-SVD approach implicitly discards $\vb_i$ and reports only $\va_i$.

    \item \textbf{Weight-independence.} Theorem~\ref{thm:weight-indep} shows that exact joint eigenfunctions ($\Xb{1}{p}\vc = \Xb{2}{p}\vc$ for all $p$) are identified by CCA regardless of the weighting, up to a small correction from the ratio of Frobenius norms across bases.
\end{enumerate}

\end{document}
